\documentclass[book.tex]{subfiles}
\begin{document}

\chapter{Imaginary Time Quantum Mechanics}
\label{chap:imaginary_qm}
\noindent\rule{11cm}{0.4pt} \\
\textbf{Prerequisites:} \autoref{chap:quantum_mechanics},   \\
\textbf{Difficulty Level:} ** \\
\noindent\rule{11cm}{0.4pt}
\vspace{5mm}
\textcolor{red}{TODO: Cite Shankar}
Quantum mechanics has built up a rich body of mathematical techniques. Can we adapt these techniques to solve problems in statistical mechanics? It turns out that a simple mathematical transformation can do the trick
\begin{align}
    t &= -i\tau
\end{align}
That is, we allow time to become a complex variable in the equations of quantum mechanics rather than a real valued variable. Applying this transformation to Schrodinger's equation, we get the imaginary-time Schrodinger equation
\begin{align}
-\hbar \frac{d|\psi(\tau)\rangle}{d\tau} &= H|\psi(\tau)\rangle
\end{align}
This operation is also commonly referred to as a "rotation" to the imaginary axis (\textcolor{red}{TODO: Insert picture of the complex plane.})
The propagator for this system is then
\begin{align}
    U(\tau) &= e^{-\frac{H}{\hbar}\tau}
\end{align}
Note that the propagator is a Hermitian operator and not a unitary operator. Let $|n\rangle$ be the eigenstates of the Hamiltonian
\begin{align}
    H|n\rangle &= E_n|n\rangle
\end{align}
We can expand the propagator in terms of these eigenstates
\begin{align}
    U(\tau) &= \sum |n\rangle \langle n | e^{-\frac{1}{\hbar}E_n\tau}
\end{align}
For the case of the harmonic oscillator, we can simply insert $t &= -i\tau$ in the closed form expression for the propagator to obtain (\textcolor{red}{TODO: Flesh out the derivation})
\begin{align}
    U(x,x';\tau) &= \frac{m\omega}{2\pi \hbar \sinh \omega \tau} \exp \left [ -\frac{m\omega}{2\hbar \sing \omega \tau}((x^2 + x'^2)\cosh \omega \tau - 2xx' \right ]
\end{align}
The Heisenberg operators $\Omega(\tau)$ can then be defined in terms of SChrodinger operator $\Omega$ (\textcolor{red}{TODO: Which Schrodinger operator is this?})
\begin{align}
    \Omega(\tau) &= e^{\frac{H}{\hbar} \tau} \Omega e^{-\frac{H}{\hbar} \tau}
\end{align}
\end{document}